% ═══════════════════════════════════════════════════════════════════
% ЧАСТЬ IX. ПОЛНЫЕ ПРОТОКОЛЫ СЕРТИФИКАТОВ
% ═══════════════════════════════════════════════════════════════════
\part{Полные протоколы сертификатов}

\section{Протокол сертификата C: перечень проверок}\label{sec:protoC}

Сертификат~C содержит полный перечень проверок эквивариантности.
Ниже --- сводная таблица групп проверок с результатами.

\begin{center}
\small
\begin{longtable}{@{}p{2.2cm} p{8.4cm} p{2.6cm}@{}}
\toprule
Группа & Проверка & Результат \\
\midrule
\endfirsthead
\toprule
Группа & Проверка & Результат \\
\midrule
\endhead
\bottomrule
\endfoot
$\mu_4$-стенд & аномалия $=i$; порядки в $\mathrm{GL}_{32}$;
элементов порядка $4$: $42$; $\mu_4$ на трёх стендах &
PASS ($4/4$) \\
Теорема C1 (K3) & изотипика $H^2$: $(1,7,7,7)$; три независимых
пути подсчёта согласованы & PASS \\
Гептада & $\Phi_7$ точен; $\Delta=-343$; $j=-3375$;
$h(-7)=1$; $\tau$ сертифицирован & PASS ($5/5$) \\
Теорема C2 (тор) & коммутатор $0$; $P^4=I$; кратность
$\lambda_1$: $4$; charpoly $x^4-1$; континуум $4\pi^2$ &
PASS ($5/5$) \\
Сетка конвейера & нарушений симметричной кодировки: $0$;
нарушений $(7,22)$: $1023/1024$; орбиты спектра согласованы &
PASS ($3/3$) \\
\midrule
\multicolumn{2}{@{}l}{Итого: $21$ проверка} & все PASS \\
\end{longtable}
\end{center}

\section{Протокол сертификата E: восемь пунктов}\label{sec:protoE}

\begin{center}
\small
\begin{longtable}{@{}p{1.2cm} p{9.6cm} p{2.4cm}@{}}
\toprule
№ & Пункт & Результат \\
\midrule
\endfirsthead
\toprule
№ & Пункт & Результат \\
\midrule
\endhead
\bottomrule
\endfoot
E-1 & Перебор $64$ пар шага: моновариант растёт, поток завершается &
PASS \\
E-2 & Конвейер $48\times48$: $t^*=48$; сверка $48/212/432/114$
полная & PASS \\
E-3 & Поправочный шаг $(7,22)$: завершимость на правилах поправок &
PASS \\
E-4 & Необратимый поток: терминальный цикл длины $4$ --- $\mu_4$-орбита
шага & PASS \\
E-5 & Масштабы $48\to384$: формула $t^*$ сохраняется & PASS \\
E-6 & K3-слой: отражения с торможением завершимы & PASS \\
E-7 & Клейн-слой: Зингер-порождение порядка $7$ & PASS \\
E-8 & Клейн-слой: $\mu_4$-орбиты длины $4$ & PASS \\
\midrule
\multicolumn{2}{@{}l}{Итого: $8$ пунктов} & $8/8$ PASS \\
\end{longtable}
\end{center}

\section{Протокол сертификата F: пять пунктов}\label{sec:protoF}

\begin{center}
\small
\begin{longtable}{@{}p{1.2cm} p{9.6cm} p{2.4cm}@{}}
\toprule
№ & Пункт & Результат \\
\midrule
\endfirsthead
\toprule
№ & Пункт & Результат \\
\midrule
\endhead
\bottomrule
\endfoot
F-1 & Единственность: $250\,000$ принятых интервалов; единственное
значение $11/4$, $q_{\min}=4$; граница $Q<7.07\cdot10^{49}$ &
PASS \\
F-2 & Априорные границы: $\det G_8=64$; K3: $Q=256$ (Крамер);
тор: $1$; Клейн: $2/1$ & PASS \\
F-3 & K3 end-to-end: $49$ измерений $\to$ восстановление точное;
макс. знаменатель $4$; $r_0-v\in\ker_{29}$ & PASS \\
F-4 & Отрицательный тест: $4\pi^2$ отвергнута (согласовано с
Линдеманом) & PASS \\
F-5 & LLL: $\im\tau=\sqrt7/2$; минполином $4X^2-7$; остаток
$7.75\cdot10^{-121}$; $j=-3375$, откл. $1.99\cdot10^{-117}$ &
PASS \\
\midrule
\multicolumn{2}{@{}l}{Итого: $5$ пунктов} & $5/5$ PASS \\
\end{longtable}
\end{center}

\section{Протокол сертификата G: шесть пунктов}\label{sec:protoG}

\begin{center}
\small
\begin{longtable}{@{}p{1.2cm} p{9.6cm} p{2.4cm}@{}}
\toprule
№ & Пункт & Результат \\
\midrule
\endfirsthead
\toprule
№ & Пункт & Результат \\
\midrule
\endhead
\bottomrule
\endfoot
G-1 & Спектр решёток: $\lambda_{m,n}=(2\pi)^2|m\tau-n|^2/(\im\tau)^2$;
$9$ решёток; семейства согласованы & PASS \\
G-2 & Дробность: приведённые $\lambda_1=4\pi^2/(\im\tau)^2$; CM ---
точные дроби $16\pi^2/d$, $4\pi^2/d$; $\pi/15$, $\pi/30$ &
PASS \\
G-3 & Поле-LLL: восстановление $(p,q,s)$ с остатком
$2.31\cdot10^{-102}$ & PASS \\
G-4 & Коула--Селберг: $d=4,7,3,15$; остатки $<10^{-100}$; $j$-пара
и $R$ & PASS \\
G-5 & K3: квадрант-интеграл $\Ga(1/4)^4/(16\pi\sqrt2)$; остаток
$3.92\cdot10^{-43}$ & PASS \\
G-6 & Лестница $\mu_N$: $2\pi/N$, $\pi/N$; $\delta_C=\pi/7$ &
PASS \\
\midrule
\multicolumn{2}{@{}l}{Итого: $6$ пунктов (время ~10 мин)} & $6/6$ PASS \\
\end{longtable}
\end{center}

\section{Протокол сертификата H: шесть пунктов}\label{sec:protoH}

\begin{center}
\small
\begin{longtable}{@{}p{1.2cm} p{9.6cm} p{2.4cm}@{}}
\toprule
№ & Пункт & Результат \\
\midrule
\endfirsthead
\toprule
№ & Пункт & Результат \\
\midrule
\endhead
\bottomrule
\endfoot
H-1 & Циклотомическая решётка: форма Рамануджа целочисленна;
SNF $(1,1,5,5,15,15,15,15)$; $\disc=1125^2$; поле
$\QQ(\zeta_{30})=\QQ(\zeta_{15})$ & PASS \\
H-2 & Объём: $\mathrm{vol}_h=(2\pi)^{32}\prod\Ga(k/N)^{-c_k}=1125$;
остаток $1.01\cdot10^{-119}$; PSLQ-дефляция ($N=15$) точна &
PASS \\
H-3 & Периоды Ферма: $\Ga$-мономы $91/406$; квадратуры
$10^{-121}$; эквивариантность точна; уровень $30$ несводим &
PASS \\
H-4 & Гросс-нормировка $d=15$: модули --- $\Ga$-мономиалы; фаза
$\psi$ восстановлена LLL; $|R|=3-\sqrt5$ & PASS \\
H-5 & Лестница $\pi/7$, $\pi/15$, $\pi/30$: веса $4\sin(\pi k/N)$
согласованы & PASS \\
H-6 & Знаменатели: $Q_{\text{стенд}}=480$; рационализация F
замыкается без скана & PASS \\
\midrule
\multicolumn{2}{@{}l}{Итого: $6$ пунктов (время ~2 мин)} & $6/6$ PASS \\
\end{longtable}
\end{center}

\section{Ключевые остатки протоколов}\label{sec:residuals}

Сводка численных остатков (максимум по каждой группе проверок):
замкнутая форма против tanh-sinh --- $3.9\cdot10^{-36}$; фазы ---
$0.0$; эквивариантность --- $3.9\cdot10^{-36}$; целостность цепей
--- $1.5\cdot10^{-36}$; лестница отражения --- $7.0\cdot10^{-36}$;
ДПФ --- $4.0\cdot10^{-35}$; глубокая запись --- $1.7\cdot10^{-71}$;
периоды сертификата~H --- $10^{-121}$; Клейн $j$ ---
$2.0\cdot10^{-117}$; слой Коула--Селберга --- $10^{-100}$;
K3-квадрант --- $3.9\cdot10^{-43}$. Все остатки согласованы с
заявленными рабочими точностями; ни один не превышает порога
сертификации.
